% FILE: 02_lineage_and_context.tex
% Project: standards
% Thesis Role: public-standards-compliance-and-gravity-field-mapping

\section{Lineage and Context}
\label{sec:standards-lineage}

\subsection{2016 Language-Model Lesson}
\label{sec:standards-lineage-2016}

The research program's foundational insight originates in a 2016 language-model experiment
that demonstrated local text imitation does not conserve consequence. A system that
predicts plausible next tokens produces text that is locally coherent but globally
unconstrained: it can assert any claim, including false process conformance claims, with
equal fluency. The experiment established that prediction without evidence is structurally
incapable of producing trustworthy operational records.

This lesson is directly relevant to the standards project. A process intelligence system
that issues conformance certificates without formal grounding in a named public standard
and a named mathematical criterion is exhibiting exactly the failure mode identified in
2016: it is generating plausible-looking output (a certificate) that has no conserved
consequence (no verifiable link to the underlying process). The standards registry is
the architectural response to this failure mode---it is the mechanism by which every
certificate is forced to name its mathematical ancestor.

The transition from prediction to coordination, which is the central theoretical
contribution of the dissertation, depends on this mechanism. Coordination requires shared
referents: two parties can coordinate on a process outcome only if they share a definition
of what counts as a valid process execution. Public standards supply those shared
referents. Without this project, the Evidence lattice is a coordination claim without
a coordination substrate.

\subsection{Enterprise Process Gap}
\label{sec:standards-lineage-enterprise-gap}

The enterprise process gap is the structural condition in which operational claims made
to boards, regulators, and transaction partners cannot be traced back to process
evidence. This gap is not a data quality problem; it is an architectural problem. The
data may exist in XES or OCEL format, but if no formal mapping connects the format to
a named validation threshold, the data cannot be used to answer the question a buyer's
advisor actually asks: ``Can you prove, to a $10^{-6}$ tolerance, that this process ran
as declared?''

The standards project closes this gap by establishing the Slide-to-Receipt bridge:

\begin{equation}
  \text{SlideBinding} = \operatorname{BLAKE3}\!\left(
    \text{SlideUUID} \parallel \text{ReceiptHash} \parallel \text{ClaimValue}
  \right)
\end{equation}

This formula, defined in \texttt{public\_standards\_to\_m\&a\_claims.md}, binds every
slide in an investment presentation containing an operational claim to a wasm4pm execution
receipt stored on the ledger. The buyer's advisors replicate the metric by executing the
same conformance query on the log in the Virtual Data Room. The standard (XES, OCEL,
Petri Net, PROV-O, or Declare) supplies the named format; the Blue River Dam threshold
($f \ge 0.95$, $p \ge 0.90$) supplies the named criterion; the BLAKE3 receipt supplies
the cryptographic binding.

\subsection{Relationship to the Chatman Equation}
\label{sec:standards-lineage-chatman-equation}

The Chatman Equation $A = \mu(O^*)$ asserts that autonomous agency is a function of
process observability over the full object lifecycle. The standards project supplies the
observability substrate: it defines, for each of the 16 public standards, what a lawful
observation looks like, what an unlawful observation looks like, and what the formal
criterion for distinguishing the two is.

Concretely, the OCEL 2.0 placement file defines valid observation as conformance to the
bipartite graph invariant: $\forall (e, o, q) \in \text{E2O},\ e \in \mathcal{E} \land
o \in \mathcal{O}$. Any event-to-object link that violates this invariant is an unlawful
observation---a dangling link refused by \texttt{OcelLog::validate} as
\texttt{OcelRefusal::DanglingEventObjectLink}. The WF-net placement file defines lawful
process completion as the soundness criteria of van der Aalst (1998): option-to-complete,
proper completion, and no dead transitions. A process that terminates without reaching the
sink place $o$ with exactly one token is an unlawful execution. These formal definitions
are the content of $\mu$: they are the measurement functions that permit the equation to
be evaluated rather than merely stated.

\subsection{Relationship to Receipts and Replay}
\label{sec:standards-lineage-receipts}

The receipt and replay infrastructure of the foundry is architecturally dependent on this
project. A receipt without a named standard is a hash without a claim. The standards
project defines what claim each hash certifies: an XES trace receipt certifies that the
trace $\sigma = \langle e_1, \ldots, e_n \rangle$ satisfies chronological invariance and
that the cumulative hash $\mathcal{H}(\sigma) = \operatorname{BLAKE3}(\mathcal{H}(e_1)
\parallel \cdots \parallel \mathcal{H}(e_n))$ is registered on the ledger. An OCEL receipt
certifies that the bipartite graph $\mathcal{L}$ passes structural validity. A WF-net
decommissioning receipt (exemplified in \texttt{wf-net\_verification\_specification.md})
records \texttt{historical\_replay\_fitness: 0.978}, \texttt{precision: 0.912}, and
\texttt{total\_recorded\_cases: 124500}---named metrics against named thresholds.

The OTel trace integration schema (\texttt{otel\_trace\_integration\_schema.json}) extends
this architecture to live telemetry: it defines the \texttt{blake3\_receipt} field per
OTel span and the \texttt{event\_chain\_root} as the BLAKE3 hash of the entire trace,
making the receipt chain continuous from the genesis span to process closure.
